\chapter{Example: the phase-driven junction}\label{ch:ex-phase}

Microwave irradiation of a superconducting loop that contains the junction modulates the
phase difference, $\phi(t)=\phi_0+A\cos\Omega t$. This is the configuration of Andreev
spectroscopy experiments, in which absorption of microwave photons excites the junction from
its ground state into the excited singlet, the \enquote{pair
transition}~\cite{bretheau2013exciting,bretheau2013supercurrent,janvier2015coherent,vanwoerkom2017microwave,hays2018direct}.
The theory of the microwave response of Andreev levels is developed in
Refs.~\cite{desposito2001controlled,zazunov2003andreev,kos2013frequencydependent,olivares2014dynamics,park2020adiabatic,kurilovich2021microwave};
for interacting dots, see
Refs.~\cite{matutecanadas2022signatures,fatemi2022microwave,hermansen2022inductive}.
Here we compute the dressed subgap states of the interacting reference junction
\eqref{eq:reference}, where $\omega_{SS^*}=E_{S^*}-E_S=\q{phase.omega_star}$.

\section{The pair transition}

For a weak drive the phase couples through the current operator
$\partial_\phi H\equiv\partial H/\partial\phi$ (\cref{sec:jacobi-anger}), and the Rabi
splitting at the resonance $\Omega\approx\omega_{SS^*}$ is
\begin{equation}\label{eq:phase-rabi}
  \delta=A\,\abs{\matel S{\partial_\phi H}{S^*}} ,
\end{equation}
with $\abs{\matel S{\partial_\phi H}{S^*}}=\q{phase.element}$ for the reference junction. The
matrix element is small: at $\eps=\q{phase.detuning}$ the junction is close to the
particle--hole-symmetric point, where it vanishes by the eta selection rule
(\cref{sec:selection,ch:ex-gate}). \Cref{fig:phase}(a) shows the rotating-frame quasienergies
$\qe_S$ and $\qe_{S^*}-\Omega$ near the resonance for $A=\q{phase.sweep_amplitude}$. The Sambe
solution in the full Fock space at QP cutoff $\q{phase.cutoff}$ and the two-level RWA give the
same avoided-crossing structure. Their energies differ by the AC Stark shift due to other states, which the
two-level RWA omits; the offset grows as $A^2$. Panel (b) compares the splitting at the
resonance with \eqref{eq:phase-rabi}; at $A=\q{phase.ratio_amplitude}$ the AC Stark shift has
lowered the resonance frequency by $\q{phase.resonance_shift}$. Both the RWA, whose coupling is
$2\abs{(H_1)_{SS^*}}=8J_1(A/4)\abs{\matel S{\partial_\phi H}{S^*}}$, and the full result fall
below the linear law as the amplitude grows: at $A=\q{phase.ratio_amplitude}$ they reach
$\q{phase.rwa_ratio}$ (RWA) and $\q{phase.sambe_ratio}$ (full Sambe result) of the linear value.
The Bessel saturation of the drive harmonic accounts for the reduction in the RWA; the
counter-rotating terms, the higher harmonics and the other states, all omitted by the
two-level RWA, reduce the effective coupling further. From
$A\approx\q{phase.hybridization_amplitude}$ onwards, the pair hybridizes with replicas of other
states and the splitting no longer follows a smooth curve.

\section{The Floquet spectrum versus frequency and amplitude}

\Cref{fig:spectrum}(a,b) shows the Floquet spectrum close to the ground state: every replica
$\qe_\alpha+\ell\Omega$ within $\pm\q{spectrum.window}$ of $E_S$, coloured by its time-averaged
weight on the undriven singlet. The calculation uses the complete Fock space of a bath with
$\q{spectrum.pairs}$ level pairs per lead (dimension $\q{spectrum.dimension}$,
$\omega_{SS^*}=\q{spectrum.omega_star}$), so that all AC Stark shifts are included. The dressed
singlet is the dark line. The light lines descending with slope $-\ell$ are replicas
$E_b-\ell\Omega$, $\ell>0$, of higher states $b$. Where a replica crosses the dressed singlet
and the drive couples the two, the lines repel and exchange weight; otherwise they cross.

As a function of the frequency, panel (a), the $S^*$ replica produces the one-photon avoided
crossing at $\Omega\approx\q{spectrum.one_photon}$ and a much narrower two-photon crossing at
$\Omega\approx\q{spectrum.two_photon}$. Above the continuum threshold,
$\Omega\gtrsim\q{spectrum.threshold}$, replicas of doublet-plus-quasiparticle states cross
densely, and the strongest of them visibly repel the dressed singlet. Away from resonances the
dressed singlet lies about $\q{spectrum.offset}\times10^{-3}$ above $E_S$. Two contributions of
comparable magnitude and opposite sign produce this shift. The time-averaged Hamiltonian $H_0$
renormalizes the hybridization, $\Gamma\to\Gamma J_0(A/4)^2$ (\cref{eq:phase-h0}); for
$A=\q{spectrum.phase_amplitude}$ the factor is $\q{spectrum.hybridization_factor}$, which weakens the
binding of the singlet and, by itself, raises its energy by $\q{spectrum.h0_shift}\times10^{-3}$. The second-order
AC Stark shifts \eqref{eq:ac-stark} lower the energy by $\q{spectrum.stark_shift}\times10^{-3}$ at
$\Omega=\q{spectrum.decomposition_omega}$, which leaves $\q{spectrum.dressed_shift}\times10^{-3}$.
At these low frequencies the result approaches the adiabatic average of the instantaneous
energy, $\langle E_S(\phi(t))\rangle-E_S=\q{spectrum.adiabatic_shift}\times10^{-3}$. As a function of the amplitude
at $\Omega=\q{spectrum.phase_omega}$, panel (b), the dressed singlet rises roughly as $A^2$.
Near $A\approx\q{spectrum.crossing}$ it hybridizes with a replica that the growing shift has
brought into resonance.

\begin{figure}[htb]
\centering
\manualfigure{spectrum}
\caption{Floquet fan diagrams of the junction ($\q{spectrum.pairs}$ level pairs per lead,
complete Fock space at QP cutoff $\q{spectrum.cutoff}$, $M=\q{spectrum.harmonics}$). Shown are all
replicas $\qe+\ell\Omega$ within $\pm\q{spectrum.window}$ of the undriven singlet energy,
coloured by their time-averaged weight on $S$. (a) Phase drive with
$A_\phi=\q{spectrum.phase_amplitude}$ versus $\Omega$; dotted lines mark the one- and
two-photon $S$--$S^*$ resonances. (b) Phase drive at $\Omega=\q{spectrum.phase_omega}$ versus
amplitude. (c, d) The same for a gate drive with $A_\eps=\q{spectrum.gate_amplitude}$ and
versus $A_\eps$.}
\label{fig:spectrum}
\end{figure}

\section{The adiabatic limit}

At low frequency, $\Omega\ll\omega_{SS^*}$, the junction follows the drive adiabatically, and
the quasienergy of the dressed ground state tends to the time average of the instantaneous
energy, \cref{eq:adiabatic},
\begin{equation}\label{eq:phase-adiabatic}
  \qe_S\simeq\frac1{2\pi}\int_0^{2\pi}E_S\bigl(\phi_0+A\cos\vartheta\bigr)\,\dd\vartheta .
\end{equation}
The Berry phase vanishes for this one-dimensional back-and-forth motion. \Cref{fig:phase}(c)
compares Floquet quasienergies at $\Omega=\q{phase.slow_omega}$ and $\q{phase.fast_omega}$ with
\eqref{eq:phase-adiabatic}. At $\Omega=\q{phase.slow_omega}$ the dressed singlet follows the
adiabatic average until the phase swing reaches $\phi=\pi$, at $A=\pi-\phi_0$, where the
$S$--$S^*$ gap is smallest. Passages through this region are no longer adiabatic. The dressed
state then acquires an admixture of $S^*$, its static weight on $S$ drops, and its quasienergy
falls below the adiabatic average:
Landau--Zener transitions at each passage interfere from cycle to cycle, the mechanism of
\cref{ch:ex-tls}~\cite{shevchenko2010landau}. At $\Omega=\q{phase.fast_omega}$ the photon
energy approaches the minimal gap; around $A\approx\q{phase.mixing_amplitude}$ a resonance with
$S^*$ mixes the two singlets strongly (faded markers: less than half of the weight on $S$), and
beyond it the dressed singlet lies above the adiabatic curve.

Since $E_S(\phi)$ is a smooth function with a minimum near $\phi=0$, the average increases with
$A$: a strong slow phase drive raises the time-averaged ground-state energy and reduces the
time-averaged supercurrent. Expanding \eqref{eq:phase-adiabatic} to second order,
$\qe_S-E_S\simeq\tfrac{A^2}4E_S''(\phi_0)$, and the curvature $E_S''$ contains the
second-order coupling to \emph{all} excited states, including quasiparticle states up to the
bandwidth. This calculation therefore uses the complete Fock space of a smaller bath; a
truncated static basis misses these contributions (\cref{ch:ex-numerics}).

\begin{figure}[htb]
\centering
\manualfigure{phase}
\caption{Phase-driven reference junction. (a) Rotating-frame quasienergies near the $S$--$S^*$
resonance for $A=\q{phase.sweep_amplitude}$, Sambe in the Fock space (solid) and RWA (dashed).
(b) Splitting at the resonance versus amplitude, with \eqref{eq:phase-rabi} and the RWA
$2\abs{(H_1)_{SS^*}}$. (c) Quasienergy of the dressed singlet at low frequencies compared with
the adiabatic average \eqref{eq:phase-adiabatic}, in the complete Fock space of a bath with
$\q{phase.adiabatic_pairs}$ level pairs per lead; faded markers denote states with less than half of their
weight on $S$.}
\label{fig:phase}
\end{figure}

\section{Continuum, multiphoton processes and experiments}

At higher frequencies the replicas of the bound states enter the continuum
(\cref{ch:ex-numerics}): for $\Omega>E_D+\Delta-E_S$ a photon can break the singlet into a
doublet plus a quasiparticle, the process that limits Andreev spectroscopy and causes
quasiparticle poisoning~\cite{zgirski2011evidence,olivares2014dynamics,riwar2015control,wesdorp2023dynamical}.
Two-photon transitions $S\to S^*$ at $\Omega\approx\omega_{SS^*}/2$ are mediated directly by the
second harmonic of the phase drive (\cref{sec:jacobi-anger}), which is proportional to the
tunnelling operator rather than to the current, and in second order by the first harmonic
through intermediate states. Floquet engineering of junctions by strong
phase or voltage drives, including steady Floquet--Andreev states, is discussed in
Refs.~\cite{peraltagavensky2018berry,venitucci2018nonadiabatic,melin2019engineering,keliri2023driven,park2022steady,ortegataberner2023anomalous}.

\begin{incode}
\begin{lstlisting}[frame=none,backgroundcolor={},aboveskip=0pt,belowskip=0pt]
static = solve(reference_model(**model), 2, Sector(0, 0), SolverOptions(method="dense", eigenpairs=60))
periodic = PeriodicHamiltonian.from_parameters(reference_model, model, ParameterDrive.cosine("phi", 1.), 0.716)
pair = solve_floquet(periodic, 2, Sector(0, 0), dict(floquet_cutoff=2, states=2), static=static)
rwa = rwa_model(periodic, [0, 1], [0, -1], 2, Sector(0, 0), static=static)
\end{lstlisting}
Passing \code{static} avoids recomputing the static eigenstates for every frequency of a scan.
\par\smallskip\usermanual{docs/floquet.md}.
\end{incode}

\tutorial{06_phase_drive}
